\documentclass[10pt]{article}
\usepackage[margin=0.75in]{geometry}
\usepackage{amsmath,amssymb,amsthm,microtype}
\newtheorem{theorem}{Theorem}
\newcommand{\R}{\mathcal R}
\newcommand{\A}{\mathcal A}
\begin{document}
\begin{center}
{\Large Auxiliary Orbits for Affine-Transfer Induction}\par\medskip
Collatz proof project\quad---\quad September 5, 2026
\end{center}
Let $T(n)=n/2$ for even $n$ and $(3n+1)/2$ for odd $n$.
Write $\R(n)$ for eventual reachability of one and
$\A(u)=[\R(3u+2)\Rightarrow\R(27u+20)]$.
The universal transfer statement is equivalent to full Collatz by the
previous formalization. We extend the available conditional induction
steps using auxiliary orbits; universal coverage remains unproved.

\begin{theorem}[Bridge rule]
Suppose natural numbers $u,v,s,t,r$ satisfy
\[
 T^s(3u+2)=3v+2,\qquad T^t(27u+20)=T^r(27v+20).
\]
Then $\A(v)\Rightarrow\A(u)$. When $v<u$, this is a valid
strong-induction step on the transfer parameter.
\end{theorem}
\begin{proof}
From $\R(3u+2)$, the first identity gives $\R(3v+2)$.
Apply $\A(v)$ and advance its larger partner $r$ steps.
The second identity then gives $\R(T^t(27u+20))$, hence
$\R(27u+20)$. Reachability is invariant under finite orbit shifts.
The rule and this invariance are kernel checked.
\end{proof}

\begin{theorem}[A progression requiring an auxiliary orbit]
For every natural $Q$, put $u=28+4096Q$ and $v=21+3072Q$.
Then $v<u$ and
\[
 T^2(3u+2)=3v+2,\qquad
 T^{12}(27u+20)=T^{10}(27v+20)=47+6561Q.
\]
Consequently $\A(21+3072Q)\Rightarrow\A(28+4096Q)$.
\end{theorem}
\begin{proof}
The base computations and odd-step counts are
\[
 T^2(86)=65\ (j=1),\quad T^{12}(776)=47\ (j=5),\quad
 T^{10}(587)=47\ (j=4).
\]
Binary-cylinder transport gives respective increments $9216Q$,
$6561Q$, and $6561Q$. Also $u-v=7+1024Q>0$.
These identities and the implication are proved for every $Q$ in
\texttt{Collatz.AffineBridge}; the smaller transfer is an explicit premise.
\end{proof}

\paragraph{Symbolic search.}
Fix a parameter cylinder $u=u_0+2^D Q$.
At clock $t\leq D$, represent a known-convergent source by
$x+3^p2^{D-t}Q$, initially $(x,p)=(3u_0+2,1)$.
An ordinary step changes $(x,p)$ to $(T(x),p+(x\bmod2))$.
If $x\equiv2\pmod3$, $(x-2)/3<u_0$, and $3^{p-1}\leq2^t$,
its associated parameter
\[
 v=(x-2)/3+3^{p-1}2^{D-t}Q
\]
is strictly smaller than $u$ for all $Q\geq0$.
Invoking $\A(v)$ permits replacement of $(x,p)$ by $(9x+2,p+2)$.
After any bounded sequence of such bridges, matching both the constant
and slope of the target certifies the original transfer conditionally.
Multiple bridges at the same clock are allowed. Search limits are explicit.

\paragraph{Finite comparison and its limits.}
For $u_0=0,\ldots,1000$, the original aligned merge and simultaneous
return criteria find 133 merges and 56 returns. The other 812 paired
orbits enter $\{1,2\}^2$ without either certificate. Equal cycle states
preserve their odd-count difference; unequal states never meet; neither
permits $y=9x+2$. Thus these 812 fail the original criteria at every later
time, not merely up to a cutoff. This is finite Python evidence, not a
Lean census theorem or an obstruction to handling finite base cases.

Allowing at most two bridges through clock 30 finds 599 conditional
certificates, including 421 of those 812; 402 parameters reach the depth
limit and none hit the state cap. Parameter 28 is among the 421.
Every saved certificate is replayed on integer lifts $Q=0,1,7,2^{32}$;
a separate small exhaustive path search checks deduplication and closure.
Sample replay does not prove universal lifting; the affine formulas above
explain it, while Lean verifies the displayed family and general bridge rule.

\paragraph{Open obligation.}
A complete proof still needs coverage by valid base cases and decreasing
recursive certificates. Auxiliary bridges enlarge the proof search, but
no global saturation or decreasing measure for unresolved states is known.
No mathematical priority claim or proof of Collatz is asserted.
\end{document}
